\begin{theorem}[The Euler sign for every degeneration and its equality criterion]
\label{cdp:st:C.2.G.1}
Let $Y$ be a connected smooth compact complex threefold with
$b_1(Y)=b_2(Y)=0$, and let $f:Y\to B$ be a proper surjective
holomorphic map with connected fibres to a smooth compact curve.
Assume that the smooth general fibre is a complex two-torus.
The full fibres may have arbitrary positive component multiplicities
and their reduced supports may have arbitrary singularities. Then
\[
 B\simeq\PP^1,\qquad
 c_3(Y)=e(Y)\in\{0,2\},\qquad b_3(Y)=2-c_3(Y).
 \tag{C.G.1}\label{cdp:eq:C.G.1}
\]
Every reduced fibre support has nonnegative Euler number.
An irreducible singular support with rational normalization
resolution has Euler number two; every other irreducible support
has Euler number zero.

The equality $c_3(Y)=0$ holds exactly when every support has one
of the following descriptions. A smooth irreducible support has
a finite unramified torus cover. A singular irreducible support
has smooth minimal elliptic ruled or smooth minimal Hopf
normalization. A reducible support has normalization a disjoint
union of smooth minimal Hopf or non-K\"ahler elliptic surfaces;
its reduced conductor image and inverse image are disjoint
unions of smooth elliptic curves, and the conductor branch maps
are finite unramified covers. Each description concerns the
normalization and resolution maps of the actual fibre. In
particular the full-fibre and conductor multiplicities are
retained.
\end{theorem}
\Uses{\ref{cdp:st:D.1.9}; \ref{cdp:st:A.3.4};
\ref{cdp:st:A.5.2}; \ref{cdp:st:A.5.3};
\ref{cdp:st:C.2.MH.1}; \ref{cdp:st:C.2.MH.2};
\ref{night:P4.26}; \ref{night:P4.32};
\ref{night:P4.36}; \ref{night:P4.37};
\ref{cdp:st:C.2.INT.1}; \ref{cdp:st:C.2.INT.2};
\ref{cdp:st:C.2.INT.3}; \ref{cdp:st:C2a};
\cite[Proposition~2.1 and Theorem~2.2, pp.~842--843]{Peters2021}}

\begin{proof}
\noindent\textbf{Step 1. The base is rational and the local
finite character is available on every full fibre.}
The beginning of connected-fibre Leray injects
$H^1(B,\Q)$ into $H^1(Y,\Q)=0$, so $B=\PP^1$.
Proposition~\ref{cdp:st:D.1.9} supplies an integral section
of $R^1f_*\Z$ which restricts nontrivially to each reduced
support and to each nearby torus. On a small disc its
Leray lift is an integral class on the entire tube.
Reduction modulo a prime which does not divide its fibre
restriction gives the finite cyclic character used in
\ref{cdp:st:C.2.MH.2}. The construction is on the original
full fibre and requires no change of its multiplicity.

\medskip\noindent\textbf{Step 2. A normal irreducible support
is smooth and has a finite unramified torus cover.}
Proposition~\ref{cdp:st:A.5.2} excludes a singular normal
support. For a smooth support $S$ in a full fibre $mS$,
the normalized change $t=s^m$ is, in a local equation
$t=uz^m$, the map
\[
 w^m=u,\qquad s=wz.
\]
These charts give a finite unramified central cover of $S$
and a proper holomorphic submersion over the new disc,
as in \ref{cdp:st:A.3.4}. Its nearby fibre is the
original torus. Smooth proper transport gives $b_1=4$
and rational first Chern class zero on the central cover.
It is K\"ahler by the even-$b_1$ criterion, and the
surface classification identifies it as a torus; the
precise classification facts are
\cite[Proposition~2.1 and Theorem~2.2, pp.~842--843]{Peters2021}.
Thus $e(S)=0$.

\medskip\noindent\textbf{Step 3. The irreducible nonnormal
normalization falls into three explicit resolution classes.}
Write $Z\to T\to S$ for the minimal resolution of the
normalization. The canonical evaluation and
\ref{cdp:st:A.5.3} give a nonempty effective divisor $D$
with $K_Z+D$ torsion. A positive pluricanonical section,
raised to a power killing this torsion, would give a
nonzero holomorphic function vanishing along a positive
multiple of $D$. Compactness excludes that function.
Hence $\kappa(Z)=-\infty$. The classification in
\cite[Theorem~2.2, p.~843]{Peters2021} gives a rational
resolution, a resolution ruled over a curve of positive
genus, or a resolution of class VII. The actual maps
$Z\to T\to S$ remain the maps in every subsequent
comparison.

\medskip\noindent\textbf{Step 4. Every irreducible support has
the claimed exact Euler value and normalization.}
For a rational resolution,
\ref{cdp:st:C.2.MH.2} gives $e(S)=2$. For a resolution
ruled over a curve of genus $g\ge1$, the same proposition
gives $e(S)=2-2g$. The local irreducible inequality in
\ref{night:P4.26} gives $e(S)\ge0$, and therefore $g=1$
and $e(S)=0$. Its proof uses the effective conductor,
the local trace inequality, and the fundamental group
of this support; no assertion about the other fibre
supports enters that local inequality. The zero-Euler
comparison in \ref{night:P4.32} then identifies $T$ as
a smooth minimal elliptic ruled surface.

For a class-VII resolution, \ref{night:P4.37} excludes
every minimal model with positive second Betti number
by its saturated-foliation Chern identity and the
negative-definite curve form. The curve-free Inoue
case cannot carry the nonempty effective conductor.
The remaining Hopf case is treated by
\ref{night:P4.36}: its derivative and foliation
calculations force the actual normalization to be
smooth and minimal, with $e(S)=0$. These arguments
are local comparisons for the given irreducible
support; they do not invoke the irreducible-or-normal-
crossings hypothesis on all fibres in
\ref{night:P4.58}. This proves the asserted exact
irreducible classification in the present scope.

\medskip\noindent\textbf{Step 5. Arbitrary reducible supports
have nonnegative Euler number, with the stated equality data.}
For each normalization resolution, an adjacent global
component supplies a nonzero effective divisor with
rationally trivial first Chern class. Its positive
pairing with a proposed K\"ahler form excludes a
K\"ahler resolution. Proposition~\ref{cdp:st:C.2.INT.1}
gives the actual Picard restriction to each elliptic
curve and excludes a Hopf--Enoki or elliptic--Enoki
interface. Proposition~\ref{cdp:st:C.2.INT.2} excludes
the remaining all-Enoki possibility using the full
multiplicities, the canonical coefficient maximum,
and the local vanishing sphere of square $-2$.

Proposition~\ref{cdp:st:C.2.INT.3} now applies to the
entire reduced support. With $A_i$ its normalization
conductor and $C$ its reduced image, its actual
normalization map gives
\[
 e(S)=\sum_i e(T_i\setminus A_i)+e(C),
 \qquad e(T_i\setminus A_i)\ge0,\quad e(C)\ge0.
\]
The proof retains the scheme conductor in the
coherent-duality argument and all contractions in
the topological calculation. It proves both
directions of the equality characterization:
the normalizations are the stated smooth minimal
surfaces, the conductor curves are smooth elliptic,
and their maps are finite unramified covers.

\medskip\noindent\textbf{Step 6. The global sum is zero or two.}
The finite-critical-value and Euler-additivity
calculation of \ref{cdp:st:C2a}, together with
complex-oriented Chern--Gauss--Bonnet, gives
\[
 c_3(Y)=e(Y)=\sum_{b\in B}e((Y_b)_{\rm red})\ge0.
 \tag{C.G.2}\label{cdp:eq:C.G.2}
\]
The smooth torus contributes zero over the regular
base. Poincar\'e duality gives $e(Y)=2-b_3(Y)$.
The nonsingular middle cup-product pairing is
skew-symmetric over $\Q$, so $b_3(Y)$ is even.
Its nonnegativity now proves \eqref{cdp:eq:C.G.1}.

\medskip\noindent\textbf{Step 7. Vanishing is equivalent to
the listed geometric data on every fibre.}
Every term in \eqref{cdp:eq:C.G.2} is nonnegative;
the sum vanishes precisely when every term does.
Steps 2 and 4 give the irreducible forms, and
Step 5 gives the reducible forms with their exact
conductor maps. Conversely each of these forms has
zero Euler number by the same local calculations,
so their occurrence on every support makes the
global sum zero.
\end{proof}
